% Journal of Bioinformatics and Computational Biology submission source.
% Uses the unmodified author-supplied World Scientific class dated 2026/05/18.
\pdfminorversion=7
\pdfsuppresswarningpagegroup=1
\documentclass{ws-jbcb}

\usepackage[super]{cite}
\usepackage{xcolor}
\usepackage[verbose,hypertexnames=false]{hyperref}
\hypersetup{colorlinks=false,allbordercolors=blue,pdfborderstyle={/S/U/W 1}}
\usepackage{array}
\usepackage{tabularx}
\usepackage{xurl}

\graphicspath{{figures/}{../../figs/}}

% theorem, proposition, definition, example, remark, and proof are class-native.
\newtheorem{algorithm}{Algorithm}
\numberwithin{equation}{section}

\newcommand{\tas}{\textit{Arabidopsis thaliana}}
\newcommand{\ath}{\textit{A.\ thaliana}}
\newcommand{\hs}{\textit{Homo sapiens}}
\newcommand{\TAU}{\tau}
\newcommand{\xRightarrow}[1]{\stackrel{#1}{\Longrightarrow}}
\newcommand{\LTS}{\mathcal{L}}
\newcommand{\powerset}{\mathcal{P}}
\newcommand{\strongsim}{\mathrel{\sim_{s}}}
\newcommand{\weaksim}{\mathrel{\sim_{w}}}
\newcommand{\weakpre}{\mathrel{\preceq_{w}}}
\newcolumntype{Y}{>{\raggedright\arraybackslash}X}

\hypersetup{
  pdftitle={A formal and reproducible framework for auditing observable behavior in qualitative biological network models},
  pdfauthor={Carlos Ramirez Ovalle},
  pdfkeywords={qualitative biological networks, Petri nets, labeled transition systems, weak bisimulation, model checking, Boolean networks}
}

\begin{document}

\markboth{C. Ramirez Ovalle}{Formal and reproducible audit of biological network models}
\catchline{}{}{}{}{}

\title{A formal and reproducible framework for auditing observable behavior in
qualitative biological network models}

\author{Carlos Ramirez Ovalle\footnote{Corresponding author.}}

\address{Department of Natural Sciences and Mathematics, Pontificia Universidad
Javeriana Cali, Calle 18 No. 118-250, Via a Pance, Cali, Valle del Cauca,
Colombia\\
carlosovalle@javerianacali.edu.co}

\maketitle

\begin{abstract}
Qualitative biological networks are often compared structurally, although
similar topology does not ensure that executable models preserve event order,
branching, or hidden steps. We present a formal, reproducible framework that maps
Petri-net and logical models to finite reachable labeled transition systems under
a declared observable interface. An implication-aware classifier reports strong
bisimilarity, weak bisimilarity, mutual or one-way weak simulation, or
non-comparability; structural and finite-trace scores remain separate
diagnostics. We prove greatest-fixed-point correctness and termination for the
implemented relation-deletion algorithm and characterize the relation hierarchy,
controlled silent refinement, and trace-equivalent branching. In six
construction-level cases, weak bisimulation classified 6/6 correctly, compared
with 5/6 for truncated traces, 4/6 for LTS-GDA, 3/6 for a structural profile, and
2/6 for PN-GDDA. Python and mCRL2 agreed on 42 strong/weak decisions, and an
independent simulation game agreed on all 67,600 ordered comparisons in the
declared small-system universe. Public GINsim models, HPN-DREAM/CASPOTS families,
and curated Petri nets test cross-format execution, held-out concordance, update
semantics, and interface dependence. The framework audits explicit models; it
does not establish predictive performance, cell-line specificity, or
organism-level equivalence.
\end{abstract}

\keywords{Qualitative biological networks; Petri nets; labeled transition
systems; weak bisimulation; model checking; Boolean networks.}

\section{Introduction}\label{sec:introduction}

Qualitative biological networks represent regulatory systems, signaling
pathways, Boolean rules, and Petri-net mechanisms when kinetic information is
limited. Models intended to describe related processes may nevertheless originate
from different organisms, inference pipelines, formalisms, or update semantics.
Their comparison is therefore a recurring computational biology problem, not
only a graph-matching task. Petri nets provide explicit causality, conflict, and
concurrency, while logical models make discrete regulatory assumptions
executable.\cite{Murata1989,ChaouiyaPetriBio2007,Li2006,Fisher2007}

Structural similarity does not determine observable executable behavior.
Graphlet methods sensitively compare local motifs; PN-GDDA, for example, uses
151 directed bipartite Petri-net graphlets and 592 published orbit
slots.\cite{Szawulak2022} By design, such a score does not encode transition
labels, initial marking, or recursive reachable branching. Conversely, internal
refinement can alter topology while preserving the behavior exposed at a chosen
interface. Structure and behavior are therefore complementary evidence layers.

Linear traces add event order but omit where alternatives become committed.
Two models can generate the same finite sequences while resolving a biological
choice at different states. Branching-time equivalences retain that information,
and simulation preorders report asymmetric containment when every behavior of one
model can be matched by another. Recent formal work on asynchronous Boolean
networks likewise highlights the role of reachability and branching structure in
phenotypic analysis.\cite{Yizengaw2026}

Let $N_1$ and $N_2$ be executable labeled models and let $\Sigma$ be a declared
observable interface. We formulate the comparison as
\begin{equation}\label{eq:comparison-map}
 (N_1,N_2,\Sigma)
 \longmapsto (\LTS(N_1),\LTS(N_2))
 \longmapsto C_{\Sigma}(\LTS(N_1),\LTS(N_2)),
\end{equation}
where $\LTS(N_i)$ is the finite reachable labeled transition system and
$C_{\Sigma}$ returns the strongest supported class among strong bisimilarity,
weak bisimilarity, mutual weak simulation, either one-way weak simulation, and
non-comparability. We call the reporting priority \emph{graded}: a stronger
relation suppresses weaker implied labels. This is an implication-aware
classifier, not an algebraic grading. Structural and truncated-trace measures
remain attached diagnostics rather than substitutes for $C_{\Sigma}$.

The formal relations are standard, but their use here is made auditable. We
define the observable interface and simulation orientation, prove
greatest-fixed-point correctness for the executed relation-deletion algorithm,
and establish the hierarchy supporting the classifier. Controlled silent
refinement and minimal counterexamples separate structure, traces, and branching.
The implementation is checked against construction-level ground truth, mCRL2,
and an independently coded attacker--defender game that is exhaustive within the
declared universe of one- and two-state systems over $\{a,\TAU\}$.

Biological evaluation then exercises different sources of heterogeneity. Public
GINsim models test SBML-qual/GINML import and asynchronous versus synchronous
updates.\cite{Naldi2018,Chaouiya2013SBML,Traynard2016,AbouJaoude2009}
HPN-DREAM/CASPOTS families connect independently inferred Boolean rules to
held-out perturbation-response data.\cite{Hill2016,Razzaq2018} Curated plant and
animal Petri nets expose dependence on the observational interface. These cases
evaluate cross-format use and empirical concordance; they are not premises of
the mathematical results.

The contributions are:
\begin{romanlist}[(iv)]
\item a formal bioinformatics framework that audits two executable qualitative
biological models under a declared interface;
\item mathematical guarantees for the implemented classifier, together with
properties and counterexamples distinguishing structure, traces, and branching;
\item independent computational validation through six controlled cases,
mCRL2, 67,600 attacker--defender comparisons, structural/trace baselines, and
observed computational scaling; and
\item applications to public GINsim models, HPN-DREAM/CASPOTS families, and
curated Petri nets, supported by public code, data provenance, tests, and an
executed notebook.
\end{romanlist}

The comparison layers retain different information. PN-GDDA and the structural
profile summarize topology; traces retain order and silent abstraction but not
branching; LTS-GDA retains labeled local motifs. The integrated audit additionally
matches recursive branching and directional containment, while held-out data
remain a separate evidence layer.

The operational sequence is summarized in Fig.~\ref{fig:pipeline}. Biological
delimitation and interface declaration precede formal encoding, reachability
construction, and classification. This separation makes the modeling
assumptions visible before any equivalence calculation is performed.

\begin{figure}[t]
\centering
\includegraphics[width=0.96\textwidth]{Fig1.pdf}
\caption{Six-stage workflow. Stages 1--3 delimit the model-comparison question
and observational interface; stages 4--6 construct executable Petri nets,
compare their reachable behavior, and assign a model-level relation.}
\label{fig:pipeline}
\alttext{A six-stage workflow from biological question and observable interface
through formal encoding, reachable-state construction, and model comparison.}
\end{figure}

\section{Formal framework for observable biological-model comparison}
\label{sec:framework}

\subsection{Labeled Petri nets and induced transition systems}
\label{sec:petri}

\begin{definition}[Labeled place/transition net]\label{def:petri-net}
A labeled place/transition net is a tuple
\begin{equation}\label{eq:petri-net}
 N=(P,T,\operatorname{pre},\operatorname{post},\ell,m_0),
\end{equation}
where $P$ and $T$ are finite sets of places and transitions,
$\operatorname{pre},\operatorname{post}:T\to\mathbb{N}^{P}$ are the input and
output multisets, $m_0\in\mathbb{N}^{P}$ is the initial marking, and
$\ell:T\to\Sigma\cup\{\TAU\}$ assigns either an observable action in the declared
interface $\Sigma$ or the internal action $\TAU$.
\end{definition}

Under standard interleaving semantics, $t$ is enabled at $m$ when
$\operatorname{pre}(t)\leq m$ componentwise; firing it changes $m$ to
$m-\operatorname{pre}(t)+\operatorname{post}(t)$. The implementation uses
a worklist to enumerate reachable markings explicitly. A state cap and a token
bound make failure to establish a finite state space explicit rather than silently
truncating the relation calculation. All curated Petri nets used in Section
\ref{sec:applications} are safe, but the definitions below apply to every finite
reachable system.

\begin{definition}[Finite labeled transition system]\label{def:lts}
A finite labeled transition system (LTS) is
\begin{equation}\label{eq:lts}
 L=(S,\Sigma_{\TAU},\longrightarrow,s_0),\qquad
 \Sigma_{\TAU}=\Sigma\cup\{\TAU\},
\end{equation}
where $S$ is a finite state set, $s_0\in S$, and
$\longrightarrow\;\subseteq S\times\Sigma_{\TAU}\times S$. We write
$s\xrightarrow{a}s'$ for $(s,a,s')\in\longrightarrow$. The reachable semantics
of the net in Definition \ref{def:petri-net} is denoted by $\LTS(N)$.
\end{definition}

\subsection{Observable and weak transitions}\label{sec:weak-transitions}

\begin{definition}[Silent closure and weak target sets]\label{def:weak-move}
For $s\in S$, let
\begin{equation}\label{eq:tau-closure}
 \operatorname{Cl}_{\TAU}(s)=
 \{u\in S:s\xrightarrow{\TAU *}u\}.
\end{equation}
For $a\in\Sigma$, define
\begin{equation}\label{eq:weak-target}
 W_a(s)=\{v:\exists u,w\quad
 s\xrightarrow{\TAU *}u\xrightarrow{a}w\xrightarrow{\TAU *}v\},
 \qquad W_{\TAU}(s)=\operatorname{Cl}_{\TAU}(s).
\end{equation}
We write $s\xRightarrow{a}v$ when $v\in W_a(s)$.
\end{definition}

The convention $W_{\TAU}=\operatorname{Cl}_{\TAU}$ is important: a direct silent
move on one side may be matched by zero or more silent moves on the other. It is
also the convention implemented by the Python engine and used in the AUT exports
checked by mCRL2.

\subsection{Simulation and bisimulation}\label{sec:relations}

Let $L_i=(S_i,\Sigma_{\TAU},\longrightarrow_i,s_i^0)$ for $i\in\{1,2\}$.

\begin{definition}[Weak simulation]\label{def:weak-simulation}
A relation $R\subseteq S_1\times S_2$ is a weak simulation from $L_1$ to $L_2$
if, whenever $(x,y)\in R$ and $x\xrightarrow{a}_1x'$, there is
$y'\in W_a^2(y)$ such that $(x',y')\in R$. We write
\begin{equation}\label{eq:simulation-orientation}
 L_1\weakpre L_2
\end{equation}
if $(s_1^0,s_2^0)$ belongs to some such relation. Thus the right-hand system
$L_2$ simulates the left-hand system $L_1$; the left-hand behavior is contained
in the right-hand behavior under the interface.
\end{definition}

\begin{definition}[Weak and strong bisimulation]\label{def:bisimulation}
A relation $R\subseteq S_1\times S_2$ is a weak bisimulation if it satisfies the
condition in Definition \ref{def:weak-simulation} and, for every $(x,y)\in R$ and
$y\xrightarrow{a}_2y'$, there is $x'\in W_a^1(x)$ with $(x',y')\in R$. The
initial systems are weakly bisimilar, written $L_1\weaksim L_2$, when their
initial pair belongs to such a relation. Strong bisimulation, written
$L_1\strongsim L_2$, replaces every weak target set $W_a^i(x)$ by the direct
target set $\{x':x\xrightarrow{a}_i x'\}$, including for $a=\TAU$.
\end{definition}

These are standard branching-time relations.\cite{Milner1989,Sangiorgi2011}
The explicit orientation in \eqref{eq:simulation-orientation} prevents the common
ambiguity between ``simulates'' and ``is simulated by'' in one-way results.

\subsection{Weak trace language}\label{sec:trace-language}

\begin{definition}[Exact and truncated weak traces]\label{def:weak-traces}
The finite weak observable language of $L$ is
\begin{equation}\label{eq:exact-language}
 L_w(L)=\{a_1\cdots a_r\in\Sigma^*:s_0
 \xRightarrow{a_1}\cdots\xRightarrow{a_r}s
 \text{ for some }s\in S\}.
\end{equation}
Its depth-$k$ restriction is
$L_{\leq k}(L)=\{u\in L_w(L):|u|\leq k\}$. The diagnostic distance is
\begin{equation}\label{eq:distance}
d_k(L_1,L_2)=1-
\frac{|L_{\leq k}(L_1)\cap L_{\leq k}(L_2)|}
{|L_{\leq k}(L_1)\cup L_{\leq k}(L_2)|}.
\end{equation}
\end{definition}

The implementation computes the prefix-closed truncated sets in Definition
\ref{def:weak-traces}. The statement $d_k=0$ is therefore not identified with
exact weak-trace equivalence. Exact weak-trace results in this study come only
from mCRL2 or from a finite construction for which the complete language is
enumerated analytically.

\subsection{Graded behavioral classification}\label{sec:decision}

\begin{definition}[Graded classifier]\label{def:classifier}
For finite $L_1,L_2$ over the declared interface $\Sigma$, define
$C_{\Sigma}(L_1,L_2)$ by the following priority:
\begin{equation}\label{eq:classifier}
C_{\Sigma}(L_1,L_2)=
\begin{cases}
\mathsf{strong}, & L_1\strongsim L_2,\\
\mathsf{weak}, & L_1\not\strongsim L_2\text{ and }L_1\weaksim L_2,\\
\mathsf{mutual}, & L_1\not\weaksim L_2,\ L_1\weakpre L_2,\ L_2\weakpre L_1,\\
\mathsf{left\_in\_right}, & L_1\weakpre L_2,\ L_2\not\weakpre L_1,\\
\mathsf{right\_in\_left}, & L_2\weakpre L_1,\ L_1\not\weakpre L_2,\\
\mathsf{noncomparable}, & \text{otherwise}.
\end{cases}
\end{equation}
The output is a reporting class, not a newly proposed equivalence relation.
\end{definition}

\begin{proposition}[Exclusive reporting]\label{prop:classifier}
The classifier in Definition \ref{def:classifier} assigns exactly one class to
every ordered pair of finite LTSs.
\end{proposition}

\begin{proof}
The first two cases are selected before any simulation-only case. If neither is
selected, the two simulation predicates have exactly four Boolean combinations:
both true, only the first true, only the second true, or both false. These are the
last four cases of \eqref{eq:classifier}, respectively. The cases are exhaustive
and disjoint by their stated guards.
\end{proof}

\subsection{Fixed-point computation and formal properties}
\label{sec:fixed-point}

For $a\in\Sigma_{\TAU}$, let $D_a^i(x)=\{x':x\xrightarrow{a}_i x'\}$ and
$U=S_1\times S_2$. Define operators on $\powerset(U)$ by
\begin{align}
\Phi_{\preceq}(R)=\{(x,y)\in U:
 &\ \forall a,x'\in D_a^1(x)\ \exists y'\in W_a^2(y),
 (x',y')\in R\},\label{eq:sim-operator}\\
\Phi_w(R)=\{(x,y)\in\Phi_{\preceq}(R):
 &\ \forall a,y'\in D_a^2(y)\ \exists x'\in W_a^1(x),
 (x',y')\in R\},\label{eq:weak-bisim-operator}\\
\Phi_s(R)=\{(x,y)\in U:
 &\ \forall a,x'\in D_a^1(x)\ \exists y'\in D_a^2(y),
 (x',y')\in R,\nonumber\\
 &\ \forall a,y'\in D_a^2(y)\ \exists x'\in D_a^1(x),
 (x',y')\in R\}.\label{eq:strong-bisim-operator}
\end{align}

\begin{theorem}[Correctness of relation deletion]\label{thm:fixed-point}
For each $\Phi\in\{\Phi_{\preceq},\Phi_w,\Phi_s\}$, the operator is monotone.
The implemented algorithm starts with $R=U$, repeatedly deletes any currently
examined pair $z\notin\Phi(R)$, and stops after a complete pass with no deletion.
On finite LTSs it terminates and returns the greatest fixed point $\nu\Phi$.
Consequently, testing membership of $(s_1^0,s_2^0)$ computes, respectively,
weak simulation in the orientation of Definition \ref{def:weak-simulation}, weak
bisimilarity, and strong bisimilarity.
\end{theorem}

\begin{proof}
If $R\subseteq Q$, every witness pair in $R$ remains available in $Q$; hence
$\Phi(R)\subseteq\Phi(Q)$ for each operator. Every deletion strictly decreases
the finite relation, so termination requires at most $|U|$ deletions.

Let $G=\nu\Phi$. The invariant $G\subseteq R$ holds initially. If it holds before
a deletion and $z\in G$, then
$z\in\Phi(G)\subseteq\Phi(R)$, so $z$ cannot be deleted. Thus $G$ remains a
subset of the terminal relation $R_*$. A final pass without deletion gives
$R_*\subseteq\Phi(R_*)$. Conversely, each $z\notin R_*$ was deleted from some
$R_z$ with $z\notin\Phi(R_z)$. Since $R_*\subseteq R_z$, monotonicity implies
$\Phi(R_*)\subseteq\Phi(R_z)$ and therefore $z\notin\Phi(R_*)$. Hence
$R_*=\Phi(R_*)$. Because $R_*$ is a fixed point, $R_*\subseteq G$; together
with the invariant this yields $R_*=G$. The three operators are exactly the
transfer clauses of Definitions \ref{def:weak-simulation} and
\ref{def:bisimulation}.
\end{proof}

\begin{remark}[Resource bound]\label{rem:complexity}
Writing $n_i=|S_i|$, the candidate relation stores $O(n_1n_2)$ pairs and permits
at most $n_1n_2$ deletions and $n_1n_2+1$ passes. Each pass scans current pairs,
outgoing moves, and matching direct or weak targets. Silent closures are cached;
visible weak targets are constructed on demand. We therefore report these
parameters instead of claiming an unjustified tight time bound.
\end{remark}

\begin{proposition}[Hierarchy of behavioral relations]\label{prop:hierarchy}
For finite LTSs under Definitions \ref{def:weak-move}--\ref{def:weak-traces},
\begin{equation}\label{eq:hierarchy}
L_1\strongsim L_2
\Longrightarrow L_1\weaksim L_2
\Longrightarrow
(L_1\weakpre L_2\ \land\ L_2\weakpre L_1)
\Longrightarrow L_w(L_1)=L_w(L_2).
\end{equation}
The converses do not hold in general.
\end{proposition}

\begin{proof}
A direct match is a permitted weak match, and each direction of a weak
bisimulation is a weak simulation. If $L_1\weakpre L_2$, induction along any
finite concrete path matches observable moves by $\TAU^*a\TAU^*$ and silent
moves by $\TAU^*$, preserving related successors. Erasing $\TAU$ gives
$L_w(L_1)\subseteq L_w(L_2)$; the reverse simulation gives equality.
Proposition \ref{prop:silent-refinement} and Example \ref{ex:branching} refute
the first and last converses. For the middle converse, compare
$a.(b+c)+a.b$ with $a.(b+c)$: each simulates the other, but the $a.b$ successor
cannot be bisimilar to a state offering both $b$ and $c$.
\end{proof}

\begin{proposition}[Controlled silent edge refinement]
\label{prop:silent-refinement}
Let $L$ contain $p\xrightarrow{a}q$ with $a\in\Sigma$ and $p\neq q$. Construct
$L'$ by copying all original states and transitions except this edge, and replacing
it by
\begin{equation}\label{eq:silent-refinement}
 \bar p\xrightarrow{a}r\xrightarrow{\TAU}\bar q,
\end{equation}
where $r$ is fresh and has no other incident behavior. Then $L\weaksim L'$.
Repeated refinements of this form also preserve weak bisimilarity. If $p$ and
$\bar p$ are initial, the selected edge is the only $a$-edge from $\bar p$ in
$L'$, and every $a$-successor of $p$ in $L$ has no outgoing $\TAU$ transition,
then $L\not\strongsim L'$.
\end{proposition}

\begin{proof}
Let
$R=\{(x,\bar x):x\in S\}\cup\{(q,r)\}$. For every copied pair other than the
source, direct moves are copied. At $(p,\bar p)$, the original edge is matched
by $\bar p\xrightarrow{a}r\xrightarrow{\TAU}\bar q$, while
$\bar p\xrightarrow{a}r$ is matched by $p\xrightarrow{a}q$ and ends in
$(q,r)$. At $(q,r)$, moves from $q$ are matched after
$r\xrightarrow{\TAU}\bar q$; that silent move is itself matched by zero silent
moves at $q$. Hence $R$ is a weak bisimulation, and repetition follows by
transitivity. Under the additional conditions, a strong match for
$\bar p\xrightarrow{a}r$ must relate $r$ to an $a$-successor $x$ of $p$, but
$r\xrightarrow{\TAU}\bar q$ would require the excluded direct $\TAU$-move from
$x$. Thus strong bisimilarity fails.
\end{proof}

\begin{proposition}[Limit of label-blind structural comparison]
\label{prop:label-blind}
Let $M$ be any comparator determined only by unlabeled structure and invariant
under relabeling of transitions. Then $M$ cannot characterize weak-trace
equality, weak simulation or weak bisimilarity for labeled systems in general.
\end{proposition}

\begin{proof}
Take two systems with the same two states and one directed edge. Label the edge
$a$ in $L_a$ and $b\neq a$ in $L_b$. Their unlabeled structures are identical,
so $M(L_a,L_b)=M(L_a,L_a)$. Yet their languages are
$\{\epsilon,a\}$ and $\{\epsilon,b\}$, and neither observable move can be
matched by the other system. Thus no decision based only on $M$ can characterize
the stated labeled relations.
\end{proof}

Proposition \ref{prop:label-blind} does not identify a defect in PN-GDDA. It
states that a label-blind structural method answers a different question and
cannot substitute for a labeled behavioral relation.

\begin{example}[Equal traces, different branching]\label{ex:branching}
Let $L_{\mathrm{late}}$ have transitions
\begin{equation}\label{eq:late-choice}
 \ell_0\xrightarrow{a}\ell_1,\qquad
 \ell_1\xrightarrow{b}\ell_b,\qquad
 \ell_1\xrightarrow{c}\ell_c,
\end{equation}
and let $L_{\mathrm{early}}$ have
\begin{equation}\label{eq:early-choice}
 e_0\xrightarrow{a}e_b,\quad e_0\xrightarrow{a}e_c,\quad
 e_b\xrightarrow{b}e_b',\quad e_c\xrightarrow{c}e_c'.
\end{equation}
Both exact prefix-closed languages are
$\{\epsilon,a,ab,ac\}$. Nevertheless,
$L_{\mathrm{late}}\not\weaksim L_{\mathrm{early}}$. There are no silent moves,
and matching either early $a$-successor with $\ell_1$ leaves one of the two
late actions unmatched. Directionally,
$L_{\mathrm{early}}\weakpre L_{\mathrm{late}}$: the late-choice state can match
the selected $b$ or $c$ continuation of either early branch. The reverse
$L_{\mathrm{late}}\weakpre L_{\mathrm{early}}$ fails because no single early
$a$-successor offers both continuations. Thus exact trace equality does not
determine branching-time equivalence or simulation orientation.
\end{example}

\subsection{Computational implementation}\label{sec:algorithms}

Algorithm~\ref{alg:pipeline} summarizes the executable path. Reachability and
the formal relations are exact whenever the finite state space is completed
below the declared state and token caps. The distance $d_k$, PN-GDDA, LTS-GDA,
and the structural profile are diagnostics rather than formal relations.

\begin{algorithm}[Executable comparison pipeline]\label{alg:pipeline}
\begin{romanlist}[(iv)]
\item Enumerate the reachable LTS from the initial marking or logical state by a
worklist; stop with an explicit error if a cap is exceeded.
\item Cache $\operatorname{Cl}_{\TAU}(s)$ and construct $W_a(s)$ by collecting
$a$-successors from the closure and the closures of their targets.
\item Initialize $R=S_1\times S_2$ and delete a pair immediately when it fails
the selected transfer operator; repeat complete passes until none is deleted.
\item Evaluate strong and weak bisimilarity and both simulation directions,
apply Eq.~(\ref{eq:classifier}), and attach trace and structural diagnostics.
\end{romanlist}
\end{algorithm}

\section{Computational implementation and validation}
\label{sec:computational-validation}

\subsection{Construction-level benchmark and comparator study}
\label{sec:benchmark-method}

Six deterministic pairs were specified before biological applications. A
12-step labeled workflow supplies identity, controlled silent subdivision, label
mismatch, label-order swap, and an added observable branch. The sixth pair is
$a.(b+c)$ versus $a.b+a.c$ from Example~\ref{ex:branching}. Table
\ref{tab:synthetic} gives construction-level and computed classes.

\begin{table}[t]
\tbl{Controlled benchmark; the expected class is not used by the classifier.
\label{tab:synthetic}}{
\begin{tabular}{@{}>{\raggedright\arraybackslash}p{0.70in}%
>{\raggedright\arraybackslash}p{1.10in}%
>{\raggedright\arraybackslash}p{1.20in}%
>{\raggedright\arraybackslash}p{1.20in}@{}}
\toprule
Case & Change & Expected & Observed \\
\colrule
Identity & exact copy & strong equivalence & strong equivalence \\
Silent refinement & insert internal steps & weak equivalence & weak equivalence \\
Label mismatch & replace an observable & not comparable & not comparable \\
Order swap & preserve counts, change order & not comparable & not comparable \\
Added branch & add an outcome & reference in candidate & reference in candidate \\
Early/late choice & move decision point & candidate in reference & candidate in reference \\
\botrule
\end{tabular}}
\end{table}

Weak bisimulation recovered all six binary weak-equivalence decisions. Truncated
trace equality recovered 5/6 but accepted the trace-equivalent branching pair;
LTS-GDA recovered 4/6, the structural profile 3/6, and direct PN-GDDA 2/6.
PN-GDDA assigned 1.0 to identity, silent refinement, label mismatch, and order
swap, while added-branch and branching pairs scored 0.9676 and 0.9657. No
threshold over the six PN-GDDA scores exceeded 4/6 because two negative controls
tied the positive controls at 1.0. Figure~\ref{fig:benchmark} reports all
predeclared decisions.

\begin{figure}[t]
\centering
\includegraphics[width=\textwidth]{Fig3.pdf}
\caption{Controlled benchmark. Left: binary decision by case and method. Right:
accuracy against labels fixed by construction. Trace equality loses branching;
PN-GDDA is label-blind; LTS-GDA retains local labels but does not abstract silent
refinement or recursively match successors.}
\label{fig:benchmark}
\alttext{A benchmark heatmap and bar chart comparing five methods on six
controlled model pairs.}
\end{figure}

The structural profile averages state-count, edge-count, label-multiset, and
out-degree-multiset similarities. The trace diagnostic declares a depth-$8$
match when $d_8=0$. PN-GDDA was implemented from its published definition using
151 directed bipartite graphlets and the 592 published/Holmes orbit
slots.\cite{Szawulak2022} For orbit $j$,
\begin{equation}\label{eq:pn-gdda}
N_G^j(k)=\frac{D_G^j(k)/k}{\sum_{r>0}D_G^j(r)/r},\quad
A^j(G,H)=1-\frac{\left[\sum_k(N_G^j(k)-N_H^j(k))^2\right]^{1/2}}{\sqrt{2}},
\end{equation}
and PN-GDDA is the mean $A^j$. Synthetic LTSs use an exact state-machine Petri
net encoding; curated pairs use their native nets. The 0.9 threshold is a fixed
diagnostic decision, not an equivalence relation. LTS-GDA instead combines four
reachable-graph orbits with labeled edge, two-step, divergence, and convergence
motifs.

The catalog audit recovered 151 graphlets and 592 slots from hash-verified Holmes
1.1.1/2.0.1.2 JARs. A direction- and type-preserving automorphism partition has
576 distinct orbits; using it is reported only as sensitivity. A Java probe
matched every Holmes 1.1.1 topology/root assignment. On a four-node test net and
its one-arc deletion, Python and Holmes scores were 0.9872027465870733 and
0.987202746587073, agreeing to 12 decimal places. Holmes is an independent
software reference for the structural score.\cite{Radom2017}

\subsection{Independent formal oracles}\label{sec:mcrl2-method}

Each analyzed LTS was exported in Aldebaran AUT format and evaluated with
\texttt{ltscompare} from hash-pinned mCRL2 202607.0.\cite{Bunte2019} The tool
independently checked strong bisimilarity, weak bisimilarity, and exact weak-trace
equivalence with \texttt{tau} hidden. Its used interface does not expose the weak
simulation preorder, so no one-way simulation is attributed to mCRL2.

A separately coded attacker--defender game grows the least set of positions from
which an unmatched move can be forced. It agreed with the production simulation
algorithm on all 67,600 ordered pairs among all 260 one- and two-state LTSs over
$\{a,\TAU\}$. This is exhaustive within that finite universe, not a proof for
arbitrary systems. Across the six synthetic pairs, six public-model controls,
and nine HPN comparisons below, Python and mCRL2 agreed on all 42 reported
strong/weak decisions. Theorem~\ref{thm:fixed-point}, finite exhaustive testing,
and external-tool comparison are distinct evidence layers.

\subsection{Observed computational scaling and reproducibility}
\label{sec:scalability-method}

Linear workflows of 8, 16, 32, 64, 96, and 128 transitions were compared with an
identity and a silently refined copy. Median runtime over three repetitions rose
from about 0.4 ms for nine states to about 0.6 s for the largest pairs. The
128-step identity relation had 16,641 candidate pairs; refinement increased it
to 18,705 pairs over 129 and 145 states. Every class remained as constructed
(the complete plot and table are archived in the repository). Timings used
Python 3.11.8 on an Apple M3 Pro with 18 GiB memory and are machine-dependent.

The Python standard library implements reachability, formal relations, parsers,
and controlled models. NumPy, pandas, Matplotlib, NetworkX, pydot, and Graphviz
support tabulation and figures. Fixed seeds, source hashes,
machine-readable outputs, unit tests, an executed notebook, environment files,
and continuous integration are public. Runtime CSV values are excluded from
byte-level regression because they depend on hardware; all scientific outcomes
and other result tables are deterministic.

\section{Biological network applications}
\label{sec:applications}

\subsection{Cross-format validation with public GINsim models}
\label{sec:public-method}

Two externally authored GINsim models test the complete import-to-comparison
path: the 13-component mammalian cell cycle in SBML-qual and the four-component
multilevel p53--Mdm2 model in GINML.\cite{Naldi2018,Traynard2016,
Chaouiya2013SBML,AbouJaoude2009} Files, URLs, and SHA-256 values are archived.
Unitary-asynchronous updates move one enabled component by one level and label
its direction. Each base model was compared with an exact copy, five silent
subdivisions, and a reachable-label perturbation. All six controls recovered
their predeclared strong/weak/exact-trace outcomes in Python and mCRL2
(Table~\ref{tab:public-controls}).

\begin{table}[t]
\tbl{Public GINsim models and control outcomes. S/W/T denote strong, weak, and
exact weak-trace equivalence.\label{tab:public-controls}}{
\begin{tabularx}{\textwidth}{@{}p{0.25\textwidth}lrrccc@{}}
\toprule
Model & Format & States & Edges & Copy & Silent & Perturbed \\
\colrule
Mammalian cell cycle & SBML-q & 826 & 3,413 & Y/Y/Y & N/Y/Y & N/N/N \\
p53--Mdm2 damage & GINML & 17 & 26 & Y/Y/Y & N/Y/Y & N/N/N \\
\botrule
\end{tabularx}}
\end{table}

Synchronous execution supplied a semantic stress test. It reached 13 states for
the cell-cycle model, all among 826 asynchronous states (Jaccard 0.0157), and
eight states for p53--Mdm2, all among 17 asynchronous states (Jaccard 0.471).
Thus the controls validate parsing, execution, export, and comparison across
formats while demonstrating that reachability is update-semantics dependent.

\subsection{HPN-DREAM/CASPOTS held-out analysis}\label{sec:hpn-method}

The public HPN-DREAM phosphoproteomic series and CASPOTS Boolean families contain
72 BT20, 191 BT549, and 21 MCF7 networks.\cite{Hill2016,Razzaq2018} One
representative per cell line was selected before test responses were read: the
model with maximum mean within-family Jaccard similarity between active clauses,
with exact lexicographic tie-breaking. This structure-only rule selected
zero-based rows 24, 137, and 9 (Table~\ref{tab:hpn-medoids}); an anti-leakage test
forbids access to test files during selection.

\begin{table}[t]
\tbl{HPN-DREAM families and structure-only representatives.
\label{tab:hpn-medoids}}{
\begin{tabular}{@{}lrrrr@{}}
\toprule
Cell line & Family & Medoid row$^a$ & Clauses & Mean Jaccard \\
\colrule
BT20  & 72  & 24  & 29 & 0.815 \\
BT549 & 191 & 137 & 25 & 0.854 \\
MCF7  & 21  & 9   & 23 & 0.851 \\
\botrule
\end{tabular}}
\begin{tabnote}$^a$Zero-based row; selected from model structure only.\end{tabnote}
\end{table}

The interface is the intersection of seven measured PI3K/MAPK/mTOR readouts.
The three common test perturbations are mTOR inhibition, IGF1 plus mTOR
inhibition, and 4EBP1 plus mTOR inhibition. Measured time-zero values initialize
readouts; stimuli and inhibitors are clamped; other relevant nodes start at zero.
Primary asynchronous LTSs label interface updates and hide all others. The nine
cross-cell pair-condition comparisons were repeated synchronously; this stress
test is not a most-permissive analysis. CASPOTS itself evaluates asynchronous
trajectories.\cite{Pauleve2020,Razzaq2018}

Experimental distance is RMSE over common post-zero times and readouts. Spearman
associations use an exact one-sided permutation of the three cell-pair labels
within each condition ($3!^3=216$). CASPOTS/Clingo compatibility scores consume
optimization to completion, use deterministic secondary tie-breaks, and were
reproduced twice. Native specificity uses the exact sign permutation of three
native-minus-mean-cross differences.

Selected LTSs had 2--44 states. None of the nine cross-cell pairs was weakly
bisimilar: six had one-way simulation and three were not comparable. The ordinal
class association with experimental distance was $\rho=0.091$ ($p=0.111$);
truncated trace distance gave $\rho=-0.596$ ($p=0.556$). LTS-GDA had the largest
observed association, $\rho=0.669$, but its exact $p=0.097$ was inconclusive.
Synchronous updating preserved seven of nine classes.

Native medoid RMSE was 0.33024 (BT20), 0.31335 (BT549), and 0.27726 (MCF7),
against discretization RMSE 0.32944, 0.30079, and 0.27726. Native excess-RMSE
advantage over mean cross-cell medoids was 0.00212 and not established by the
exact sign test ($p=0.25$; Fig.~\ref{fig:hpn-validation}). The result demonstrates
held-out compatibility and separates formal similarity from empirical response
similarity; it does not demonstrate cell-line predictive discrimination.

\begin{figure}[t]
\centering
\includegraphics[width=\textwidth]{Fig8.pdf}
\caption{Held-out HPN-DREAM analysis. (a) CASPOTS excess RMSE by source medoid
and target; boxes mark native pairs. (b) Between-cell profile RMSE by formal
class. Exact tests are condition-stratified; neither panel establishes native
cell-line specificity.}
\label{fig:hpn-validation}
\alttext{CASPOTS excess-error matrix and boxplots relating formal classes to
held-out response distances for three breast-cancer cell lines.}
\end{figure}

\subsection{Curated Petri nets and interface dependence}\label{sec:case-method}

Five plant--animal pairs model DNA repair/genome instability (GIM), energy
metabolism (DCE), cell-cycle control (SPS), cell death/autophagy (RCD), and
immune control (AID). Comparative studies and cancer-hallmark literature guided
component selection; dedicated reviews informed plant DNA-damage, cell-cycle,
and autophagy mechanisms.\cite{Quimbaya2012,ClavijoBuritica2023,Hanahan2022,
Yoshiyama2013} Transition-level
provenance is machine-checked. These manually curated qualitative nets contain
no kinetics, concentrations, tissue context, or independent experimental
validation.

Figure~\ref{fig:gim-petri} exposes the executable GIM encodings. Both include
damage sensing, checkpoint control, repair, and restoration. The animal model
uses apoptosis, whereas the plant model uses an internal SMR step followed by
differentiation or endoreduplication. Their shared \texttt{exit\_cycle} label is
a deliberately coarse observation, not evidence that terminal mechanisms are
biologically identical.

\begin{figure}[t]
\centering
\begin{minipage}[t]{\textwidth}
  \centering\includegraphics[width=0.98\textwidth]{Fig2a.pdf}
  \par\smallskip\textbf{(a) Animal/human GIM model}
\end{minipage}
\par\smallskip
\begin{minipage}[t]{\textwidth}
  \centering\includegraphics[width=0.98\textwidth]{Fig2b.pdf}
  \par\smallskip\textbf{(b) Arabidopsis GIM model}
\end{minipage}
\caption{Curated GIM Petri nets: (a) animal/human and (b) Arabidopsis. Circles
are places, rounded boxes are transitions, brackets identify observables, and
gray $\tau$ transitions are internal. The diagrams and analysis use the same
executable definitions.}
\label{fig:gim-petri}
\alttext{Two Petri-net diagrams for animal and Arabidopsis DNA-damage modules,
showing places, observable transitions, and internal transitions.}
\end{figure}

The ten nets produced 5--14 reachable states. GIM, DCE, and SPS were weakly but
not strongly bisimilar; RCD gave plant-in-animal one-way simulation; AID was
non-comparable. Table~\ref{tab:case} juxtaposes these classes with native
PN-GDDA. The 592- and 576-orbit variants changed no score by 0.001 and no
threshold decision: all five pairs were structurally similar at 0.9 despite
three different behavioral classes.

\begin{table}[t]
\tbl{Curated model-pair outputs. Relations and $d_6$ are interface-conditional;
PN-GDDA is structural.\label{tab:case}}{
\begin{tabular}{@{}l>{\raggedright\arraybackslash}p{1.05in}>{\raggedright\arraybackslash}p{1.05in}rrr@{}}
\toprule
Module & Process & Formal class & $d_6$ & PN-592 & Orbit-576 \\
\colrule
GIM & DNA repair & weak bisimulation & 0.000 & 0.9976 & 0.9975 \\
DCE & energy metabolism & weak bisimulation & 0.000 & 0.9969 & 0.9968 \\
SPS & cell-cycle control & weak bisimulation & 0.000 & 1.0000 & 1.0000 \\
RCD & death/autophagy & plant in animal & 0.143 & 1.0000 & 1.0000 \\
AID & immune control & non-comparable & 0.429 & 0.9907 & 0.9904 \\
\botrule
\end{tabular}}
\end{table}

The GIM reachable systems make the weak relation inspectable in the repository:
damage, checkpoint, repair, restoration, and cycle-exit paths match, while
additional internal paths are absorbed by $\TAU$ closure. A finer interface that
distinguishes apoptosis, differentiation, and endoreduplication would alter the
classification.

RCD supplies the contrasting directional result: both nets contain stress,
mTOR inhibition, autophagy, and survival, but only the animal encoding exposes a
canonical caspase-apoptosis branch; the plant branch uses an internal metacaspase
event. Silent subdivision preserved every class over 300 refinements per module,
while targeted relabeling changed the expected relations. Distances against
2,000 label permutations assess sensitivity inside the curated models, not
external biological significance. Full RCD diagrams and robustness plots remain
in the reproducible repository outputs.

\section{Discussion}\label{sec:discussion}

\subsection{What agreement means}
\label{sec:discussion-meaning}

The classifier turns an ordered pair of executable models into one statement at
a declared observational interface. Proposition \ref{prop:hierarchy} supports
the reporting priority: strong bisimilarity entails weak bisimilarity, and weak
bisimilarity entails both simulation directions. One-way classes retain their
orientation and therefore express behavioral containment, not equivalence. The
fixed-point theorem concerns the algorithm actually executed, including its
in-place deletion order and treatment of $\TAU$; software tests complement that
proof but cannot replace it.

The controlled cases also show why structural, linear-time, and branching
comparisons answer different questions. Identical unlabeled structure cannot
determine labeled behavior (Proposition \ref{prop:label-blind}), while Example
\ref{ex:branching} has identical exact finite traces but different branching.
PN-GDDA scored 2/6 on the predeclared behavioral task, although the independent
Holmes implementation reproduced every structural score to 12 decimal places.
This is not a defect in PN-GDDA: labels, initial states, and reachable branching
are outside that statistic by design. LTS-GDA recovered labeled local motifs and
scored 4/6, but it still lacked silent abstraction and recursive successor
matching. Its strongest held-out association ($\rho=0.669$) remained
inconclusive ($p=0.097$). The evidence therefore supports complementary reporting
of structural and behavioral measures, not replacement of one by the other.

\subsection{Interface and update semantics}
\label{sec:discussion-interface}

Every relation is indexed by $\Sigma$. Declaring an event internal changes the
paths hidden by abstraction, while merging outcomes under one label coarsens the
question. In the GIM comparison, apoptosis and differentiation or
endoreduplication share \texttt{exit\_cycle}; weak bisimilarity consequently
means agreement only at that resolution, not identity of the terminal biological
mechanisms. Proposition \ref{prop:silent-refinement} covers the controlled
transition subdivisions used here, but not arbitrary refinements that introduce
conflicts, interleavings, or observable choices.

Update semantics is likewise part of the comparison contract. Switching between
asynchronous and synchronous execution changed public GINsim behavior and two HPN
classes. General-asynchronous, most-permissive, and stochastic conventions were
not evaluated.\cite{Pauleve2020} A reported relation should therefore always name
both its observable interface and execution semantics.

\subsection{Evidence and biological limits}
\label{sec:discussion-biological}

Python and mCRL2 agreed on all 42 strong/weak decisions across the synthetic,
public-model, and HPN experiments. The independent attacker--defender checker
also agreed on all 67,600 ordered pairs in the complete declared universe of 260
one- and two-state LTSs over $\{a,\TAU\}$. These results are evidence against
small-system implementation defects, not a proof beyond the theorem. The GINsim
transformations test ingestion and execution but do not constitute independent
biological labels; the distinct HPN model families and held-out measurements add
an empirical check.

That check remains deliberately modest. The HPN analysis contains nine
pair--condition observations: the formal-class association with experimental
distance was $\rho=0.091$ ($p=0.111$), truncated trace distance was also
unsupported, and the CASPOTS native-versus-cross advantage was not established
($p=0.25$). CASPOTS RMSE measures minimum-distance compatibility over admissible
trajectories, not a unique time-series prediction.\cite{Razzaq2018} These data do
not support prognostic or cell-line-specific claims.

The curated Petri nets omit kinetics, quantities, noise, tissue context, and
independent experimental inference. Provenance makes modeling choices auditable
but does not turn them into independent evidence. Label-permutation results
measure sensitivity inside those constructions, and cross-organism weak
relations do not imply organism-level functional equivalence.

\subsection{Computational scope and use}
\label{sec:discussion-computation}

The implementation explicitly enumerates states and may store
$|S_1||S_2|$ candidate pairs. The scaling experiment reached 145 states and
18,705 pairs, so it does not establish genome-scale performance. Symbolic state
representations, partition refinement, and compositional checking are natural
extensions. PN-GDDA was applied to native curated nets and canonical
state-machine encodings; using it on SBML-qual or Boolean rules would require a
non-unique reverse translation. Larger repository models, prospectively
expert-labeled relations, biological replicates, blinded interface selection,
and additional update semantics are required before broader predictive claims.
At present, the framework is best used to audit two explicit qualitative models
at a predeclared observational resolution.

\section{Conclusion}\label{sec:conclusion}

Comparing executable qualitative biological models requires an explicit account
of observables, silent events, initial states, branching, and update semantics.
The proposed classifier reports the strongest supported relation from strong
bisimilarity through one-way weak simulation, with fixed-point correctness and
termination established for finite LTSs. Controlled cases, mCRL2, and an
independent simulation game verify distinct parts of the implementation, while
the direct PN-GDDA comparison identifies the information unavailable to a purely
structural score. GINsim, HPN-DREAM/CASPOTS, and curated Petri-net applications
show that the audit can cross model formats without hiding interface-dependent
assumptions. The present evidence supports reproducible model comparison, not
prognosis, cell-line specificity, or organism-level equivalence.

\section*{Data and code availability}

Code, curated definitions, hash-verified public models, normalized response
files, metadata, and generated results are available under the MIT License at
\url{https://github.com/cero1979/BisimulationMol}. The repository documents
validation, figure regeneration, tests, and notebook execution. No confidential
or participant-level data were used.

\section*{Declarations}

\noindent\textbf{Funding.} No specific funding was received.
\textbf{Competing interests.} The author declares no competing interests.
\textbf{Ethics.} Only computational models and published aggregate data were
used; no human participants, samples, or live animals were involved.
\textbf{Author contributions.} Carlos Ramirez Ovalle performed the
conceptualization, methodology, software, formal analysis, validation,
visualization, drafting, and revision.

\section*{Disclosure of generative AI assistance}

OpenAI Codex assisted with language editing, organization, journal-format
adaptation, and reproducibility checks. The author reviewed the manuscript and
is responsible for its content, results, interpretations, and conclusions.

\bibliographystyle{ws-jbcb}
\bibliography{references_jbcb}

\end{document}
